\chapter{Pneumococcal disease}
\index{Pneumococcal disease}

\medskip
\noindent\textit{Educational information; seek professional advice for individual care.}

\section*{What it is}
Pneumococcal disease is caused by the bacterium \textit{Streptococcus pneumoniae}. It can cause less serious infections such as ear or sinus infection, or invasive disease such as pneumonia, meningitis, or bloodstream infection.

\section*{Symptoms}
Symptoms depend on the site of infection. Pneumonia may cause fever, chills, cough, chest pain, fast breathing, or breathlessness. Meningitis may cause severe headache, neck stiffness, vomiting, confusion, light sensitivity, or a seizure. Babies, older adults, and people with weakened immunity may have less typical signs.

\section*{Diagnosis and treatment}
Assessment may include examination, oxygen measurement, chest imaging, blood tests, and samples from blood, sputum, or spinal fluid when invasive disease is suspected. Antibiotics are used when bacterial disease is diagnosed or strongly suspected; severe infection may require hospital care and support for breathing or circulation.

\section*{Prevention}
Vaccination is an important prevention measure, with the recommended product and schedule depending on age, health conditions, pregnancy, and local policy. Hand hygiene, respiratory etiquette, avoiding tobacco smoke, and managing chronic illness also reduce risk. Ask a clinician which vaccines are appropriate.

\section*{When to seek urgent care}
Seek emergency help for difficulty breathing, blue or grey lips, confusion, a seizure, severe headache with neck stiffness, collapse, or rapidly worsening fever or illness.

\section*{Sources}
\begin{itemize}
\item World Health Organization: \url{https://www.who.int/teams/health-product-policy-and-standards/standards-and-specifications/norms-and-standards/vaccine-standardization/pneumococcal-disease}
\end{itemize}
